\section{The regular-cardinal branch}

\import{function.tex}
\import{combinatorics/graph-cardinality.tex}
\import{cardinal.tex}
\import{cardinal/comparison.tex}
\import{cardinal/cofinal-families.tex}
\import{cardinal/cofinality.tex}
\import{recursion/transfinite.tex}
\import{set/choice.tex}

\subsection{Forbidden points from a history}

\begin{definition}\label{edm_history_neighbourhoods}
    $\edmHistoryNeighbourhoods{E}{X}{g}
        =\{(x,\graphNeighbourhood{E}{x}{X})\mid x\in\ran{g}\}$.
\end{definition}

\begin{proposition}\label{edm_history_neighbourhoods_iff}
    $(x,A)\in\edmHistoryNeighbourhoods{E}{X}{g}$ iff
        $x\in\ran{g}$ and $A=\graphNeighbourhood{E}{x}{X}$.
\end{proposition}
\begin{proof}
    Follows by \auto{edm_history_neighbourhoods,pair_eq_iff}.
\end{proof}

\begin{theorem}\label{edm_history_neighbourhoods_type}
    $\edmHistoryNeighbourhoods{E}{X}{g}
        \in\funs{\ran{g}}{\pow{X}}$.
\end{theorem}
\begin{proof}
    Let $N=\edmHistoryNeighbourhoods{E}{X}{g}$.
    $N$ is a relation by
        \auto{edm_history_neighbourhoods,relation}.
    $N$ is right-unique by
        \auto{edm_history_neighbourhoods,rightunique,pair_eq_iff}.
    $\dom{N}=\ran{g}$ by
        \auto{setext,dom_iff,edm_history_neighbourhoods,pair_eq_iff}.
    For every $x\in\ran{g}$ we have
        $\graphNeighbourhood{E}{x}{X}\in\pow{X}$ by
        \obv{graph_neighbourhood_subseteq,powerset_intro}.
    For every $x\in\dom{N}$ we have
        $N(x)=\graphNeighbourhood{E}{x}{X}$ by
        \auto{edm_history_neighbourhoods,function_apply_intro}.
    Hence $N$ is a function to $\pow{X}$ by \auto.
    Follows by \auto{funs_intro}.
\end{proof}

\begin{proposition}\label{edm_history_neighbourhoods_apply}
    Let $x\in\ran{g}$.
    Then $(\edmHistoryNeighbourhoods{E}{X}{g})(x)
        =\graphNeighbourhood{E}{x}{X}$.
\end{proposition}
\begin{proof}
    $\edmHistoryNeighbourhoods{E}{X}{g}$ is a function by
        \auto{edm_history_neighbourhoods_type,funs_is_function}.
    $(x,\graphNeighbourhood{E}{x}{X})
        \in\edmHistoryNeighbourhoods{E}{X}{g}$ by
        \obv{edm_history_neighbourhoods}.
    Follows by \obv{function_apply_intro}.
\end{proof}

\begin{definition}\label{edm_forbidden_set}
    $\edmForbiddenSet{E}{X}{g}
        =\ran{g}\union
            \unions{\ran{\edmHistoryNeighbourhoods{E}{X}{g}}}$.
\end{definition}

\begin{proposition}\label{edm_forbidden_set_iff}
    $z\in\edmForbiddenSet{E}{X}{g}$ iff
        $z\in\ran{g}$ or there exists $x\in\ran{g}$ such that
            $z\in\graphNeighbourhood{E}{x}{X}$.
\end{proposition}
\begin{proof}
    Let $N=\edmHistoryNeighbourhoods{E}{X}{g}$.
    $N\in\funs{\ran{g}}{\pow{X}}$ by
        \auto{edm_history_neighbourhoods_type}.
    Hence $N$ is a function and $\dom{N}=\ran{g}$ by
        \obv{funs_is_function,funs}.
    Show that if $z\in\edmForbiddenSet{E}{X}{g}$, then
        $z\in\ran{g}$ or there exists $x\in\ran{g}$ such that
            $z\in\graphNeighbourhood{E}{x}{X}$.
    \begin{subproof}
        Assume $z\in\edmForbiddenSet{E}{X}{g}$.
        Hence $z\in\ran{g}$ or
            $z\in\unions{\ran{N}}$ by
            \auto{edm_forbidden_set,union_iff}.
        \begin{byCase}
            \caseOf{$z\in\ran{g}$.}
                Follows.
            \caseOf{$z\notin\ran{g}$.}
                Hence $z\in\unions{\ran{N}}$.
                Take $A\in\ran{N}$ such that $z\in A$ by
                    \obv{unions_iff}.
                Take $x\in\dom{N}$ such that $N(x)=A$ by
                    \obv{fun_ran_iff}.
                Hence $x\in\ran{g}$.
                $N(x)=\graphNeighbourhood{E}{x}{X}$ by
                    \auto{edm_history_neighbourhoods_apply}.
                Hence $z\in\graphNeighbourhood{E}{x}{X}$.
        \end{byCase}
    \end{subproof}
    Show that if $z\in\ran{g}$ or there exists $x\in\ran{g}$ such that
            $z\in\graphNeighbourhood{E}{x}{X}$, then
        $z\in\edmForbiddenSet{E}{X}{g}$.
    \begin{subproof}
        Assume $z\in\ran{g}$ or there exists $x\in\ran{g}$ such that
            $z\in\graphNeighbourhood{E}{x}{X}$.
        \begin{byCase}
            \caseOf{$z\in\ran{g}$.}
                Follows by \auto{edm_forbidden_set,union_iff}.
            \caseOf{$z\notin\ran{g}$.}
                Take $x\in\ran{g}$ such that
                    $z\in\graphNeighbourhood{E}{x}{X}$.
                $N(x)=\graphNeighbourhood{E}{x}{X}$ by
                    \auto{edm_history_neighbourhoods_apply}.
                $N(x)\in\ran{N}$ by \auto{fun_ran_iff}.
                Hence $z\in\unions{\ran{N}}$ by \auto{unions_iff}.
                Follows by \auto{edm_forbidden_set,union_iff}.
        \end{byCase}
    \end{subproof}
    Follows by \auto.
\end{proof}

\begin{proposition}\label{edm_history_neighbourhood_union_subseteq}
    $\unions{\ran{\edmHistoryNeighbourhoods{E}{X}{g}}}
        \subseteq X$.
\end{proposition}
\begin{proof}
    $\edmHistoryNeighbourhoods{E}{X}{g}
        \in\funs{\ran{g}}{\pow{X}}$ by
        \obv{edm_history_neighbourhoods_type}.
    Hence every element of
        $\ran{\edmHistoryNeighbourhoods{E}{X}{g}}$ is a subset of $X$ by
        \auto{funs_ran,subseteq,powerset_elim}.
    Follows by \auto{subseteq,unions_iff}.
\end{proof}

\begin{proposition}\label{edm_forbidden_set_subseteq}
    Suppose $\ran{g}\subseteq X$.
    Then $\edmForbiddenSet{E}{X}{g}\subseteq X$.
\end{proposition}
\begin{proof}
    $\unions{\ran{\edmHistoryNeighbourhoods{E}{X}{g}}}
        \subseteq X$ by
        \obv{edm_history_neighbourhood_union_subseteq}.
    Follows by \auto{edm_forbidden_set,subseteq_union_iff}.
\end{proof}

\subsection{Availability at a regular-cardinal stage}

\begin{theorem}\label{edm_regular_stage_available}
    Let $\kappa$ be a regular cardinal.
    Suppose $X\approx\kappa$.
    Let $\beta\in\kappa$.
    Let $g\in\funs{\beta}{X}$.
    Suppose for every $x\in X$ we have
        $\cardLt{\graphNeighbourhood{E}{x}{X}}{\kappa}$.
    Then $\cardLt{\edmForbiddenSet{E}{X}{g}}{\kappa}$ and
        $X\setminus\edmForbiddenSet{E}{X}{g}\approx\kappa$.
\end{theorem}
\begin{proof}
    $g$ is a function and $\dom{g}=\beta$ by
        \obv{funs_is_function,funs}.
    Hence $\cardLe{\ran{g}}{\dom{g}}$ by
        \obv{function_range_cardinal_le_domain}.
    Hence $\cardLe{\ran{g}}{\beta}$.
    $\cardLt{\beta}{\kappa}$ by \obv{cardinal_element_small}.
    Hence $\cardLt{\ran{g}}{\kappa}$ by
        \obv{cardinal_le_lt_transitive}.
    $\ran{g}\subseteq X$ by \obv{funs_ran}.
    Let $N=\edmHistoryNeighbourhoods{E}{X}{g}$.
    $N\in\funs{\ran{g}}{\pow{X}}$ by
        \auto{edm_history_neighbourhoods_type}.
    Hence $N$ is a function and $\dom{N}=\ran{g}$ by
        \obv{funs_is_function,funs}.
    For every $x\in\ran{g}$ we have
        $N(x)=\graphNeighbourhood{E}{x}{X}$ by
        \auto{edm_history_neighbourhoods_apply}.
    Hence for every $x\in\ran{g}$ we have
        $\cardLt{N(x)}{\kappa}$ by \auto{subseteq}.
    Hence $\cardLt{\unions{\ran{N}}}{\kappa}$ by
        \auto{regular_indexed_union}.
    Hence $\cardLt{\ran{g}\union\unions{\ran{N}}}{\kappa}$ by
        \obv{small_binary_union}.
    Hence $\cardLt{\edmForbiddenSet{E}{X}{g}}{\kappa}$ by
        \auto{edm_forbidden_set}.
    Hence $X\setminus\edmForbiddenSet{E}{X}{g}\approx\kappa$ by
        \obv{large_minus_small}.
\end{proof}

\subsection{The total greedy step}

\begin{definition}\label{edm_greedy_step}
    $\edmGreedyStep{E}{X}{\kappa}
        =\{w\in
            ((\kappa\times\pow{\kappa\times X})\times X)\mid
            \text{
            ($X\setminus
                \edmForbiddenSet{E}{X}{\snd{\fst{w}}}$ is inhabited and
                $\snd{w}=\choice{X\setminus
                    \edmForbiddenSet{E}{X}{\snd{\fst{w}}}}$) or
            ($X\setminus
                \edmForbiddenSet{E}{X}{\snd{\fst{w}}}=\emptyset$ and
                $\snd{w}=\choice{X}$)}\}$.
\end{definition}

\begin{proposition}\label{edm_greedy_step_iff}
    Let $q\in\kappa\times\pow{\kappa\times X}$ and $y\in X$.
    Then $(q,y)\in\edmGreedyStep{E}{X}{\kappa}$ iff
        ($X\setminus\edmForbiddenSet{E}{X}{\snd{q}}$ is inhabited and
            $y=\choice{X\setminus
                \edmForbiddenSet{E}{X}{\snd{q}}}$) or
        ($X\setminus\edmForbiddenSet{E}{X}{\snd{q}}=\emptyset$ and
            $y=\choice{X}$).
\end{proposition}
\begin{proof}
    $(q,y)\in
        (\kappa\times\pow{\kappa\times X})\times X$ by
        \obv{times_tuple_intro}.
    Follows by \auto{edm_greedy_step,fst_eq,snd_eq}.
\end{proof}

\begin{theorem}\label{edm_greedy_step_type}
    Suppose $X$ is inhabited.
    Then $\edmGreedyStep{E}{X}{\kappa}
        \in\funs{\kappa\times\pow{\kappa\times X}}{X}$.
\end{theorem}
\begin{proof}
    Let $D=\kappa\times\pow{\kappa\times X}$.
    Let $T=\edmGreedyStep{E}{X}{\kappa}$.
    $T\subseteq D\times X$ by
        \auto{edm_greedy_step,subseteq}.
    Hence $T$ is a relation by
        \auto{relation,subseteq,times_elem_is_tuple}.
    Show that $T$ is right-unique.
    \begin{subproof}
        Show that for all $q,y,z$ such that $(q,y),(q,z)\in T$
            we have $y=z$.
        \begin{subproof}
            Fix $q,y,z$.
            Assume $(q,y),(q,z)\in T$.
            Hence $q\in D$ and $y,z\in X$ by
                \obv{subseteq,times_tuple_elim}.
            Hence
                ($X\setminus\edmForbiddenSet{E}{X}{\snd{q}}$
                    is inhabited and
                    $y=\choice{X\setminus
                        \edmForbiddenSet{E}{X}{\snd{q}}}$) or
                ($X\setminus\edmForbiddenSet{E}{X}{\snd{q}}=\emptyset$ and
                    $y=\choice{X}$) by
                \auto{edm_greedy_step_iff}.
            Hence
                ($X\setminus\edmForbiddenSet{E}{X}{\snd{q}}$
                    is inhabited and
                    $z=\choice{X\setminus
                        \edmForbiddenSet{E}{X}{\snd{q}}}$) or
                ($X\setminus\edmForbiddenSet{E}{X}{\snd{q}}=\emptyset$ and
                    $z=\choice{X}$) by
                \auto{edm_greedy_step_iff}.
            Follows by \auto{emptyset}.
        \end{subproof}
        Follows by \obv{rightunique}.
    \end{subproof}
    Show that $\dom{T}=D$.
    \begin{subproof}
        Show that $\dom{T}\subseteq D$.
        \begin{subproof}
            Follows by \auto{subseteq,dom_iff,times_tuple_elim}.
        \end{subproof}
        Show that $D\subseteq\dom{T}$.
        \begin{subproof}
            Show that every element of $D$ is an element of $\dom{T}$.
            \begin{subproof}
                Fix $q\in D$.
                $X\setminus\edmForbiddenSet{E}{X}{\snd{q}}\subseteq X$ by
                    \obv{setminus_subseteq}.
                Let $A=X\setminus
                    \edmForbiddenSet{E}{X}{\snd{q}}$.
                $A\subseteq X$.
                \begin{byCase}
                    \caseOf{$A$ is inhabited.}
                        Hence $\choice{A}\in A$ by \obv{choice_in}.
                        Hence $\choice{A}\in X$ by \obv{subseteq}.
                        $(q,\choice{A})\in T$ by
                            \auto{edm_greedy_step_iff}.
                        Follows by \auto{dom_iff}.
                    \caseOf{$A$ is not inhabited.}
                        Hence $A=\emptyset$ by
                            \obv{nonempty_inhabited}.
                        $\choice{X}\in X$ by \obv{choice_in}.
                        $(q,\choice{X})\in T$ by
                            \auto{edm_greedy_step_iff}.
                        Follows by \auto{dom_iff}.
                \end{byCase}
            \end{subproof}
            Follows by \obv{subseteq}.
        \end{subproof}
        Follows by \obv{subseteq_antisymmetric}.
    \end{subproof}
    Show that for every $q\in\dom{T}$ we have $T(q)\in X$.
    \begin{subproof}
        Fix $q\in\dom{T}$.
        $(q,T(q))\in T$ by \obv{function_apply_elim}.
        Hence $(q,T(q))\in D\times X$ by \obv{subseteq}.
        Follows by \obv{times_tuple_elim}.
    \end{subproof}
    Hence $T$ is a function to $X$.
    Follows by \auto{funs_intro}.
\end{proof}

\begin{proposition}\label{edm_greedy_step_available}
    Suppose $X$ is inhabited.
    Let $\beta\in\kappa$ and
        $g\in\pow{\kappa\times X}$.
    Suppose $X\setminus\edmForbiddenSet{E}{X}{g}$ is inhabited.
    Then $(\edmGreedyStep{E}{X}{\kappa})((\beta,g))
        =\choice{X\setminus\edmForbiddenSet{E}{X}{g}}$ and
        $(\edmGreedyStep{E}{X}{\kappa})((\beta,g))
            \in X\setminus\edmForbiddenSet{E}{X}{g}$.
\end{proposition}
\begin{proof}
    $(\beta,g)\in\kappa\times\pow{\kappa\times X}$ by
        \obv{times_tuple_intro}.
    Let $q=(\beta,g)$.
    $q\in\kappa\times\pow{\kappa\times X}$.
    $\edmGreedyStep{E}{X}{\kappa}
        \in\funs{\kappa\times\pow{\kappa\times X}}{X}$ by
        \obv{edm_greedy_step_type}.
    Hence $\edmGreedyStep{E}{X}{\kappa}$ is a function by
        \obv{funs_is_function}.
    $\choice{X\setminus\edmForbiddenSet{E}{X}{g}}
        \in X\setminus\edmForbiddenSet{E}{X}{g}$ by \obv{choice_in}.
    Hence $\choice{X\setminus\edmForbiddenSet{E}{X}{g}}\in X$ by
        \obv{setminus_subseteq,subseteq}.
    $(q,\choice{X\setminus\edmForbiddenSet{E}{X}{g}})
        \in\edmGreedyStep{E}{X}{\kappa}$ by
        \auto{edm_greedy_step_iff,fst_eq,snd_eq}.
    Hence $(\edmGreedyStep{E}{X}{\kappa})(q)
        =\choice{X\setminus\edmForbiddenSet{E}{X}{g}}$ by
        \obv{function_apply_intro}.
\end{proof}

\begin{proposition}\label{edm_greedy_step_empty}
    Suppose $X$ is inhabited.
    Let $\beta\in\kappa$ and
        $g\in\pow{\kappa\times X}$.
    Suppose $X\setminus\edmForbiddenSet{E}{X}{g}=\emptyset$.
    Then $(\edmGreedyStep{E}{X}{\kappa})((\beta,g))=\choice{X}$.
\end{proposition}
\begin{proof}
    $(\beta,g)\in\kappa\times\pow{\kappa\times X}$ by
        \obv{times_tuple_intro}.
    Let $q=(\beta,g)$.
    $q\in\kappa\times\pow{\kappa\times X}$.
    $\edmGreedyStep{E}{X}{\kappa}
        \in\funs{\kappa\times\pow{\kappa\times X}}{X}$ by
        \obv{edm_greedy_step_type}.
    Hence $\edmGreedyStep{E}{X}{\kappa}$ is a function by
        \obv{funs_is_function}.
    $\choice{X}\in X$ by \obv{choice_in}.
    Hence $(q,\choice{X})\in\edmGreedyStep{E}{X}{\kappa}$ by
        \auto{edm_greedy_step_iff,fst_eq,snd_eq}.
    Hence $(\edmGreedyStep{E}{X}{\kappa})(q)=\choice{X}$ by
        \obv{function_apply_intro}.
\end{proof}

\subsection{The greedy recursion}

\begin{proposition}\label{edm_greedy_recursion_type}
    Let $\kappa$ be an ordinal.
    Suppose $X$ is inhabited.
    Then
    \[
        \ordinalRecursion{\kappa}{X}{\edmGreedyStep{E}{X}{\kappa}}
            \in\funs{\kappa}{X}
    \].
\end{proposition}
\begin{proof}
    $\edmGreedyStep{E}{X}{\kappa}
        \in\funs{\kappa\times\pow{\kappa\times X}}{X}$ by
        \obv{edm_greedy_step_type}.
    Follows by \obv{ordinal_recursion_type}.
\end{proof}

\begin{proposition}\label{edm_greedy_recursion_equation}
    Let $\kappa$ be an ordinal.
    Suppose $X$ is inhabited.
    Let $\beta\in\kappa$.
    Then
    \[
        (\ordinalRecursion{\kappa}{X}
            {\edmGreedyStep{E}{X}{\kappa}})(\beta)
        =(\edmGreedyStep{E}{X}{\kappa})
            ((\beta,\restrl{\ordinalRecursion{\kappa}{X}
                {\edmGreedyStep{E}{X}{\kappa}}}{\beta}))
    \].
\end{proposition}
\begin{proof}
    $\edmGreedyStep{E}{X}{\kappa}
        \in\funs{\kappa\times\pow{\kappa\times X}}{X}$ by
        \obv{edm_greedy_step_type}.
    Follows by \obv{ordinal_recursion_equation}.
\end{proof}

\begin{theorem}\label{edm_greedy_new_point}
    Let $\kappa$ be a regular cardinal.
    Suppose $X\approx\kappa$.
    Let $\beta\in\kappa$.
    Suppose for every $x\in X$ we have
        $\cardLt{\graphNeighbourhood{E}{x}{X}}{\kappa}$.
    Then
    \[
        (\ordinalRecursion{\kappa}{X}
            {\edmGreedyStep{E}{X}{\kappa}})(\beta)
        \in X\setminus\edmForbiddenSet{E}{X}
            {\restrl{\ordinalRecursion{\kappa}{X}
                {\edmGreedyStep{E}{X}{\kappa}}}{\beta}}
    \].
\end{theorem}
\begin{proof}
    $\kappa$ is infinite by definition.
    Hence $\kappa$ is an ordinal by \auto{cardinal,initial_ordinal}.
    $\kappa$ is inhabited by \obv{infinite_is_inhabited}.
    Hence $X$ is inhabited by \obv{equinum_inhabited}.
    $\ordinalRecursion{\kappa}{X}{\edmGreedyStep{E}{X}{\kappa}}
        \in\funs{\kappa}{X}$ by
        \obv{edm_greedy_recursion_type}.
    $(\ordinalRecursion{\kappa}{X}
            {\edmGreedyStep{E}{X}{\kappa}})(\beta)
        =(\edmGreedyStep{E}{X}{\kappa})
            ((\beta,\restrl{\ordinalRecursion{\kappa}{X}
                {\edmGreedyStep{E}{X}{\kappa}}}{\beta}))$ by
        \obv{edm_greedy_recursion_equation}.
    Let $h=\ordinalRecursion{\kappa}{X}
        {\edmGreedyStep{E}{X}{\kappa}}$.
    $h\in\funs{\kappa}{X}$ and
        $h(\beta)=(\edmGreedyStep{E}{X}{\kappa})
            ((\beta,\restrl{h}{\beta}))$.
    $\beta$ is an ordinal by \obv{elem_of_ordinal_is_ordinal}.
    Hence $\beta\subseteq\kappa$ by
        \obv{ordinal_elem_implies_subseteq}.
    Hence $\restrl{h}{\beta}\in\funs{\beta}{X}$ by
        \obv{funs_restrl}.
    $h\in\rels{\kappa}{X}$ by \obv{funs}.
    Hence $h\in\pow{\kappa\times X}$ by \obv{rels}.
    Hence $h\subseteq\kappa\times X$ by \obv{powerset_elim}.
    $\restrl{h}{\beta}\subseteq h$ by \obv{restrl_subseteq}.
    Hence $\restrl{h}{\beta}\subseteq\kappa\times X$ by
        \obv{subseteq_transitive}.
    Hence $\restrl{h}{\beta}\in\pow{\kappa\times X}$ by
        \obv{powerset_intro}.
    $X\setminus\edmForbiddenSet{E}{X}{\restrl{h}{\beta}}
        \approx\kappa$ by \auto{edm_regular_stage_available}.
    Hence $X\setminus\edmForbiddenSet{E}{X}{\restrl{h}{\beta}}$
        is inhabited by \obv{equinum_inhabited}.
    Hence $(\edmGreedyStep{E}{X}{\kappa})
            ((\beta,\restrl{h}{\beta}))
        \in X\setminus
            \edmForbiddenSet{E}{X}{\restrl{h}{\beta}}$ by
        \auto{edm_greedy_step_available}.
    Hence $h(\beta)\in
        X\setminus\edmForbiddenSet{E}{X}{\restrl{h}{\beta}}$.
\end{proof}

\begin{theorem}\label{edm_greedy_later_separated}
    Let $\kappa$ be a regular cardinal.
    Suppose $X\approx\kappa$.
    Let $\beta\in\kappa$ and $\alpha\in\beta$.
    Suppose for every $x\in X$ we have
        $\cardLt{\graphNeighbourhood{E}{x}{X}}{\kappa}$.
    Then
    \[
        (\ordinalRecursion{\kappa}{X}
            {\edmGreedyStep{E}{X}{\kappa}})(\alpha)
        \neq
        (\ordinalRecursion{\kappa}{X}
            {\edmGreedyStep{E}{X}{\kappa}})(\beta)
    \]
    and
    \[
        (\ordinalRecursion{\kappa}{X}
            {\edmGreedyStep{E}{X}{\kappa}})(\alpha)
        \not\mathrel{E}
        (\ordinalRecursion{\kappa}{X}
            {\edmGreedyStep{E}{X}{\kappa}})(\beta)
    \].
\end{theorem}
\begin{proof}
    Hence $\kappa$ is a cardinal and $\kappa$ is infinite by definition.
    Hence $\kappa$ is an ordinal by \auto{cardinal,initial_ordinal}.
    $\kappa$ is inhabited by \obv{infinite_is_inhabited}.
    Hence $X$ is inhabited by \obv{equinum_inhabited}.
    $\ordinalRecursion{\kappa}{X}{\edmGreedyStep{E}{X}{\kappa}}
        \in\funs{\kappa}{X}$ by
        \obv{edm_greedy_recursion_type}.
    $(\ordinalRecursion{\kappa}{X}
            {\edmGreedyStep{E}{X}{\kappa}})(\beta)
        \in X\setminus\edmForbiddenSet{E}{X}
            {\restrl{\ordinalRecursion{\kappa}{X}
                {\edmGreedyStep{E}{X}{\kappa}}}{\beta}}$ by
        \obv{edm_greedy_new_point}.
    Let $h=\ordinalRecursion{\kappa}{X}
        {\edmGreedyStep{E}{X}{\kappa}}$.
    $h\in\funs{\kappa}{X}$ and
        $h(\beta)\in X\setminus
            \edmForbiddenSet{E}{X}{\restrl{h}{\beta}}$.
    Hence $h$ is a function and $\dom{h}=\kappa$ by
        \obv{funs_is_function,funs}.
    $\beta$ is an ordinal by \obv{elem_of_ordinal_is_ordinal}.
    Hence $\beta\subseteq\kappa$ by
        \obv{ordinal_elem_implies_subseteq}.
    Hence $\beta\subseteq\dom{h}$.
    $\restrl{h}{\beta}\in\funs{\beta}{X}$ by
        \obv{funs_restrl}.
    Hence $\restrl{h}{\beta}$ is a function and
        $\dom{\restrl{h}{\beta}}=\beta$ by
        \obv{funs_is_function,funs}.
    $(\restrl{h}{\beta})(\alpha)=h(\alpha)$ by
        \obv{restrl_of_function_apply}.
    Hence $h(\alpha)\in\ran{\restrl{h}{\beta}}$ by
        \auto{fun_ran_iff}.
    Hence $h(\alpha)\in
        \edmForbiddenSet{E}{X}{\restrl{h}{\beta}}$ by
        \auto{edm_forbidden_set_iff}.
    $h(\alpha),h(\beta)\in X$ by \obv{funs_type_apply}.
    $h(\beta)\notin
        \edmForbiddenSet{E}{X}{\restrl{h}{\beta}}$ by
        \obv{setminus}.
    Hence $h(\alpha)\neq h(\beta)$.
    Show that $h(\alpha)\not\mathrel{E}h(\beta)$.
    \begin{subproof}
        Suppose $h(\alpha)\mathrel{E}h(\beta)$.
        Hence $h(\beta)\in
            \graphNeighbourhood{E}{h(\alpha)}{X}$ by
            \auto{graph_neighbourhood_iff}.
        Hence $h(\beta)\in
            \edmForbiddenSet{E}{X}{\restrl{h}{\beta}}$ by
            \auto{edm_forbidden_set_iff}.
        Contradiction.
    \end{subproof}
\end{proof}

\begin{theorem}\label{regular_small_degree_independent}
    Let $\kappa$ be a regular cardinal.
    Suppose $X\approx\kappa$.
    Suppose $E$ is a simple graph on $X$.
    Suppose for every $x\in X$ we have
        $\cardLt{\graphNeighbourhood{E}{x}{X}}{\kappa}$.
    Then $E$ has an independent set in $X$ equinumerous with $\kappa$.
\end{theorem}
\begin{proof}
    Hence $\kappa$ is a cardinal and $\kappa$ is infinite by definition.
    Hence $\kappa$ is an ordinal by \auto{cardinal,initial_ordinal}.
    $\kappa$ is inhabited by \obv{infinite_is_inhabited}.
    Hence $X$ is inhabited by \obv{equinum_inhabited}.
    $\ordinalRecursion{\kappa}{X}{\edmGreedyStep{E}{X}{\kappa}}
        \in\funs{\kappa}{X}$ by
        \obv{edm_greedy_recursion_type}.
    Let $h=\ordinalRecursion{\kappa}{X}
        {\edmGreedyStep{E}{X}{\kappa}}$.
    $h\in\funs{\kappa}{X}$.
    Hence $h$ is a function and $\dom{h}=\kappa$ by
        \obv{funs_is_function,funs}.
    Show that $h$ is injective.
    \begin{subproof}
        Show that for all $\alpha,\beta\in\dom{h}$ such that
            $h(\alpha)=h(\beta)$ we have $\alpha=\beta$.
        \begin{subproof}
            Fix $\alpha,\beta\in\dom{h}$.
            Assume $h(\alpha)=h(\beta)$.
            Hence $\alpha,\beta\in\kappa$.
            \begin{byCase}
                \caseOf{$\alpha=\beta$.}
                    Follows.
                \caseOf{$\alpha\neq\beta$.}
                    $\alpha$ and $\beta$ are ordinals by
                        \obv{elem_of_ordinal_is_ordinal}.
                    $\alpha\in\beta$ or $\beta\in\alpha$ by
                        \obv{ordinal_elem_connex}.
                    \begin{byCase}
                        \caseOf{$\alpha\in\beta$.}
                            $h(\alpha)\neq h(\beta)$ by
                                \auto{edm_greedy_later_separated}.
                            Contradiction.
                        \caseOf{$\alpha\notin\beta$.}
                            Hence $\beta\in\alpha$.
                            $h(\beta)\neq h(\alpha)$ by
                                \auto{edm_greedy_later_separated}.
                            Contradiction.
                    \end{byCase}
            \end{byCase}
        \end{subproof}
        Follows by \obv{injective_function}.
    \end{subproof}
    Hence $h\in\bijections{\dom{h}}{\ran{h}}$ by
        \obv{injective_function_bijection_to_range}.
    Hence $h\in\bijections{\kappa}{\ran{h}}$.
    Let $H=\ran{h}$.
    $\kappa\approx H$ by \obv{equinum}.
    Hence $H\approx\kappa$ by \obv{equinum_sym}.
    $H\subseteq X$ by \auto{funs_ran}.
    Show that for all $x,y\in H$ such that $x\neq y$ we have
        $x\not\mathrel{E}y$.
    \begin{subproof}
        Fix $x,y\in H$.
        Assume $x\neq y$.
        Take $\alpha\in\dom{h}$ such that $h(\alpha)=x$ by
            \obv{fun_ran_iff}.
        Take $\beta\in\dom{h}$ such that $h(\beta)=y$ by
            \obv{fun_ran_iff}.
        Hence $\alpha,\beta\in\kappa$ and $\alpha\neq\beta$.
        $\alpha$ and $\beta$ are ordinals by
            \obv{elem_of_ordinal_is_ordinal}.
        $\alpha\in\beta$ or $\beta\in\alpha$ by
            \obv{ordinal_elem_connex}.
        \begin{byCase}
            \caseOf{$\alpha\in\beta$.}
                $h(\alpha)\not\mathrel{E}h(\beta)$ by
                    \auto{edm_greedy_later_separated}.
            \caseOf{$\alpha\notin\beta$.}
                Hence $\beta\in\alpha$.
                $h(\beta)\not\mathrel{E}h(\alpha)$ by
                    \auto{edm_greedy_later_separated}.
                Show that $h(\alpha)\not\mathrel{E}h(\beta)$.
                \begin{subproof}
                    Suppose $h(\alpha)\mathrel{E}h(\beta)$.
                    Hence $h(\beta)\mathrel{E}h(\alpha)$ by
                        \obv{simple_graph_edge_symmetric}.
                    Contradiction.
                \end{subproof}
        \end{byCase}
    \end{subproof}
    Hence $H$ is an independent set in $X$ with respect to $E$ by
        \obv{independent_in_graph}.
\end{proof}
